%
% claude-session-end-unordered-owned.tex  (FROZEN NEGATIVE CONTROL: do not edit)
%
% NewHello is enabled while the old end is in progress, but ReleaseTty keeps the owner check.  Expected: the invariant HOLDS.
%
% Z Specification for how a Claude session ends when biffd owns its record:
% EndOwnSession (a connection of the current run ends its own session),
% ExpireClaude (retention removes a session nobody is attached to), the
% ordering of a clear (/clear ends the old session before the new hello), and the
% obligation that no row outlives an end plus retention.
%
% Companion to session-model.tex (SessionEnd, RetentionExpire for CLI sessions,
% the daemon's records and runs) and to sentinel-reap.tex (BadOrphan, findings G1
% and G2 of the sentinel mechanism).  It is kept apart because those two
% specifications already have large state spaces, and this one needs the effects
% of an end (the relay row, the tty-name reservation, the logout event, the held
% mail), which neither carries.
%
% Type-check: fuzz -t claude-session-end.tex
% Model-check: see the end of this file.
%

\documentclass[a4paper,10pt,fleqn]{article}
\usepackage[margin=1in]{geometry}
\usepackage{fuzz}
\usepackage[colorlinks=true,linkcolor=blue,citecolor=blue,urlcolor=blue]{hyperref}

\begin{document}

\title{Biff Claude Session End: A Z Specification}
\author{Punt Labs}
\date{November 2026}
\maketitle

%%----------------------------------------------------------------------
\section{Introduction}
%%----------------------------------------------------------------------

Until biffd owned the session record, a Claude session ended in its own process.
The process wrote a sentinel, tried to delete its row and release its tty name,
and a reaper in some other process finished what the first left.
\texttt{sentinel-reap.tex} models that mechanism and finds that a row can be left
with no process and no sentinel behind it ($BadOrphan$): a crash in one step (G1),
and the fallback sentinel racing its own reaper in five (G2).

With the record in the daemon the question changes.  The record is meant to be
durable (the code keeps it on disk, below) and
the daemon, not a process that is going away, performs the end.  Two transitions
end a Claude session:

\begin{itemize}
  \item $EndOwnSession$: a connection of the \emph{current} run of the session
    ends its own session.  The daemon releases the tty name, deletes the row,
    appends a logout event and files the mail it holds for the session back into
    the user's inbox.  A connection of a run that a newer run has replaced is
    refused and changes nothing.
  \item $ExpireClaude$: a session with no connection that has been dormant for the
    retention period is removed with the same four effects.
\end{itemize}

Each of the four effects is its own step, and the record stays until all four are
done.  The obligations are:

\begin{enumerate}
  \item \textbf{No row outlives its record.}  The relay row and the tty name are
    held only while the record exists or is being ended, so there is no
    $BadOrphan$: the daemon has no state in which a row has nothing behind it.
  \item \textbf{No row outlives an end plus retention.}  A session nobody is
    attached to is eventually removed, once its retention is spent.
  \item \textbf{No mail is lost by an end.}  Mail issued is held by the session,
    still in the user's inbox, or shown; an end moves what it holds back to the
    inbox.
  \item \textbf{A clear ends the old session before the new hello.}  The new
    session's tty reservation is never released by the old session's end.
  \item \textbf{A replaced run cannot end the session.}
\end{enumerate}

%%----------------------------------------------------------------------
\section{Basic Types}
%%----------------------------------------------------------------------

\begin{zed}
  ZBOOL ::= ztrue | zfalse
\end{zed}

\begin{zed}
  RecState ::= rsNone | rsLive | rsDormant | rsEnding | rsStopped
\end{zed}

\begin{zed}
  TtyOwner ::= tNone | tOld | tNew
\end{zed}

\begin{zed}
  DropCause ::= dNone | dRepoGone | dRowGone | dFailedStarts | dForgot
\end{zed}

\begin{zed}
  StopKind ::= sNone | sGraceful | sTimedOut | sProcess
\end{zed}

\begin{axdef}
  maxAge : \nat \\
  maxMail : \nat \\
  maxFails : \nat
  \where
  maxAge = 2 \\
  maxMail = 2 \\
  maxFails = 2
\end{axdef}

$rsNone$ is no record (never started, or ended), $rsLive$ a record with a
connection, $rsDormant$ a record without one, and $rsEnding$ a record whose end
has begun and not finished.  $maxAge$ is the retention, in sweep ticks.
$tOld$ and $tNew$ are the old and the new session of a clear; both use the same
tty name, which is the reason the order matters.  Runs are $0$ and $1$.

%%----------------------------------------------------------------------
\section{State}
%%----------------------------------------------------------------------

\begin{schema}{Record}
  recState : RecState \\
  currentRun : 0 \upto 1 \\
  conns : \power (0 \upto 1) \\
  age : 0 \upto maxAge \\
  endRun : 0 \upto 1
\end{schema}

The daemon's record: its state, the current run of the session, the runs that
have a connection open (the connection of a replaced run is still open, and
refused), and the sweep ticks since the last connection left, and (a ghost) the run whose
connection began the end in progress.

\begin{schema}{Effects}
  rowPresent : ZBOOL \\
  ttyOwner : TtyOwner \\
  logged : ZBOOL \\
  rowTodo : ZBOOL \\
  ttyTodo : ZBOOL \\
  logTodo : ZBOOL \\
  mailTodo : ZBOOL \\
  newLive : ZBOOL \\
  ownEnded : ZBOOL
\end{schema}

What an end touches: the relay row, the tty-name reservation, the logout event,
and the four obligations an end in progress still owes.  $newLive$ is the new
session of a clear, after its hello.

\begin{schema}{Mail}
  outboxN : 0 \upto maxMail \\
  sessN : 0 \upto maxMail \\
  userN : 0 \upto maxMail \\
  shownN : 0 \upto maxMail \\
  issuedN : 0 \upto maxMail \\
  lostN : 0 \upto maxMail \\
  strandedN : 0 \upto maxMail \\
  flightN : 0 \upto maxMail \\
  shownDormant : ZBOOL \\
  forgot : ZBOOL \\
  lastStop : StopKind \\
  stopHeld : 0 \upto maxMail \\
  stopLost : 0 \upto maxMail \\
  lostBefore : 0 \upto maxMail \\
  repoGone : ZBOOL \\
  starts : 0 \upto maxFails - 1 \\
  dropCause : DropCause
\end{schema}

Mail counts: held by the session ($outboxN$), in the session's own inbox in the
relay ($sessN$), in the user's mailbox ($userN$, where an end files what it
drains),
shown to the client ($shownN$), lost with a daemon that restarted ($lostN$, a
ghost), and issued in all ($issuedN$, a ghost).  $strandedN$ counts the mail left in the relay inbox of a session whose
record was dropped ($sessN$ at that moment): nothing drains or reads it.
Mail already in the user's mailbox ($userN$) is not touched by the loss of a
record.  $shownDormant$ is a ghost
that records a print made while nobody was connected.  $forgot$ is a ghost set by
$DropRecord$: the record is gone while the row and the name may remain.  $dropCause$ is the cause of the
last drop ($dForgot$ is used only by the control that does not keep records).  $repoGone$ is the
environment: the repository directory is gone.  $starts$ counts the starts in a row that could not
build the session again.
$flightN$ is mail a read has taken from the relay and not yet handed over to the
outbox: the consuming fetch has run and the read has not finished.  $lastStop$ is
a ghost that records how the daemon last stopped: $sGraceful$ (everything given
back), $sTimedOut$ (a graceful stop whose bounded wait ended with a read still in
flight), $sProcess$ (the process was lost).  $stopHeld$ and $stopLost$ are ghosts of that last
stop alone: the mail the stop found held ($outboxN + flightN$ at the moment of the stop) and the
part of it the stop lost.  Neither is cumulative, so a goal on them cannot be met by mail an
earlier crash lost.  $lostBefore$ is $lostN$ as the last stop found it, so that
$lostN = lostBefore + stopLost$ ties what a stop reports losing to what it added to
$lostN$: only a stop changes $lostN$, so a stop that adds dropped mail to $lostN$
without reporting it in $stopLost$ breaks the equation.

\begin{schema}{EndState}
  Record \\
  Effects \\
  Mail
\end{schema}

\begin{schema}{EndInv}
  EndState
  \where
  rowPresent = ztrue \implies ( recState \neq rsNone \lor forgot = ztrue ) \\
  ttyOwner = tOld \implies ( recState \neq rsNone \lor forgot = ztrue ) \\
  newLive = ztrue \implies ttyOwner = tNew \\
  recState = rsDormant \implies conns = \emptyset \\
  recState = rsStopped \implies
    ( conns = \emptyset \land outboxN = 0 \land flightN = 0 ) \\
  recState = rsLive \implies conns \neq \emptyset \\
  recState = rsEnding \implies endRun = currentRun \\
  recState \neq rsEnding \implies
    ( rowTodo = zfalse \land ttyTodo = zfalse \land logTodo = zfalse \land
      mailTodo = zfalse ) \\
  outboxN + flightN + sessN + userN + shownN + lostN + strandedN = issuedN \\
  lastStop = sGraceful \implies stopLost = 0 \\
  stopLost \leq stopHeld \\
  lostN = lostBefore + stopLost \\
  strandedN > 0 \implies
    dropCause \in \{ dRepoGone, dRowGone, dFailedStarts \} \\
  shownDormant = zfalse
\end{schema}

The invariants, in the order of the obligations.  A row, or the old session's tty
name, exists only while the record does ($BadOrphan$ is not a state) unless the
record was dropped ($forgot$, below).  The new
session of a clear always holds its tty name.  A live record has a connection and
a dormant one has none.  An end that is not in progress owes nothing.  And mail
is conserved: issued is held, in the session's inbox, in the user's mailbox,
shown, lost with a stop of the daemon (the outbox is its memory), or stranded
by a drop of the record, the only way mail is stranded.  An end loses nothing: $FileOutbox$ files
what is held and drains the session's inbox into the user's mailbox.  A stop that
settled ($sGraceful$) lost nothing of its own ($stopLost = 0$); a stop never loses
more than it found held, and what it reports losing is what $lostN$ gained.

\begin{schema}{Invariant}
  EndInv
\end{schema}

\begin{schema}{Init}
  EndState'
  \where
  recState' = rsLive \\
  currentRun' = 0 \\
  conns' = \{ 0 \} \\
  age' = 0 \\
  endRun' = 0 \\
  rowPresent' = ztrue \\
  ttyOwner' = tOld \\
  logged' = zfalse \\
  rowTodo' = zfalse \\
  ttyTodo' = zfalse \\
  logTodo' = zfalse \\
  mailTodo' = zfalse \\
  newLive' = zfalse \\
  ownEnded' = zfalse \\
  outboxN' = 0 \\
  sessN' = 0 \\
  userN' = 0 \\
  shownN' = 0 \\
  issuedN' = 0 \\
  lostN' = 0 \\
  strandedN' = 0 \\
  flightN' = 0 \\
  lastStop' = sNone \\
  stopHeld' = 0 \\
  stopLost' = 0 \\
  lostBefore' = 0 \\
  repoGone' = zfalse \\
  starts' = 0 \\
  dropCause' = dNone \\
  shownDormant' = zfalse \\
  forgot' = zfalse
\end{schema}

%%----------------------------------------------------------------------
\section{Operations on the record}
%%----------------------------------------------------------------------

\begin{schema}{ReplaceRun}
  \Delta Record \\
  \Xi Effects \\
  \Xi Mail
  \where
  recState = rsLive \\
  currentRun = 0 \\
  recState' = rsLive \\
  currentRun' = 1 \\
  conns' = conns \cup \{ 1 \} \\
  age' = age \\
  endRun' = endRun
\end{schema}

A resume with a fresh agent opens run $1$.  The connection of run $0$ stays open;
it belongs to a replaced run.

\begin{schema}{Detach}
  \Delta Record \\
  \Xi Effects \\
  \Xi Mail \\
  run? : 0 \upto 1
  \where
  recState = rsLive \\
  run? \in conns \\
  conns' = conns \setminus \{ run? \} \\
  recState' = ( \IF conns' = \emptyset \THEN rsDormant \ELSE rsLive ) \\
  currentRun' = currentRun \\
  age' = 0 \\
  endRun' = endRun
\end{schema}

A connection goes.  Detaching ends nothing: the record becomes dormant when its
last connection goes.

\begin{schema}{Reattach}
  \Delta Record \\
  \Xi Effects \\
  \Xi Mail
  \where
  recState = rsDormant \\
  recState' = rsLive \\
  currentRun' = currentRun \\
  conns' = \{ currentRun \} \\
  age' = 0 \\
  endRun' = endRun
\end{schema}

\begin{schema}{SweepTick}
  \Delta Record \\
  \Xi Effects \\
  \Xi Mail
  \where
  recState = rsDormant \\
  age < maxAge \\
  recState' = recState \\
  currentRun' = currentRun \\
  conns' = conns \\
  age' = age + 1 \\
  endRun' = endRun
\end{schema}

The daemon's sweep counts a dormant record's age.



A stop leaves the record on disk, and the mail differs between the two stops: $GracefulStop$ waits up to the
request timeout for the reads in flight ($settled? = ztrue$ when they finish in time), files every
session's held outbox back to the relay (\texttt{Handback.close\_relays}), and closes the relays; if
the wait times out ($settled? = zfalse$, $sTimedOut$) the mail of a read still between its consuming
fetch and its hand-over is lost with the process, the one loss a graceful stop has.  $ProcessLoss$ (a
crash or a kill) loses the outbox and the mail in flight ($lostN$, PL-1753).  The mail waiting in the
relay is lost by neither, and a restored session finds it where it was.

\begin{schema}{RowAgesOut}
  \Xi Record \\
  \Delta Effects \\
  \Xi Mail
  \where
  ( recState = rsStopped \lor ( forgot = ztrue \land recState = rsNone ) ) \\
  rowPresent = ztrue \\
  rowPresent' = zfalse \\
  ttyOwner' = ttyOwner \\
  logged' = logged \\
  rowTodo' = rowTodo \\
  ttyTodo' = ttyTodo \\
  logTodo' = logTodo \\
  mailTodo' = mailTodo \\
  newLive' = newLive \\
  ownEnded' = ownEnded
\end{schema}

\begin{schema}{TtyAgesOut}
  \Xi Record \\
  \Delta Effects \\
  \Xi Mail
  \where
  forgot = ztrue \\
  recState = rsNone \\
  ttyOwner = tOld \\
  ttyOwner' = tNone \\
  rowPresent' = rowPresent \\
  logged' = logged \\
  rowTodo' = rowTodo \\
  ttyTodo' = ttyTodo \\
  logTodo' = logTodo \\
  mailTodo' = mailTodo \\
  newLive' = newLive \\
  ownEnded' = ownEnded
\end{schema}

The relay's row lifetime and the name reservation's lifetime run out, while the daemon is down or after
a drop.  This is the
NATS relay.  The local relay (no relay URL configured) keeps a reservation as a
lockfile with no lifetime: there $TtyAgesOut$ does not happen, and the session of a dropped
record holds its name until something releases it.  The model has no local
variant: the NATS behaviour is the favourable one.

\begin{schema}{GracefulStop}
  \Delta Record \\
  \Xi Effects \\
  \Delta Mail \\
  settled? : ZBOOL
  \where
  recState \in \{ rsLive, rsDormant \} \\
  recState' = rsStopped \\
  currentRun' = currentRun \\
  conns' = \emptyset \\
  age' = ( \IF recState = rsLive \THEN 0 \ELSE age ) \\
  endRun' = endRun \\
  outboxN' = 0 \\
  flightN' = 0 \\
  sessN' = sessN + outboxN + ( \IF settled? = ztrue \THEN flightN \ELSE 0 ) \\
  lostN' = lostN + ( \IF settled? = ztrue \THEN 0 \ELSE flightN ) \\
  userN' = userN \\
  shownN' = shownN \\
  issuedN' = issuedN \\
  shownDormant' = shownDormant \\
  strandedN' = strandedN \\
  forgot' = forgot \\
  repoGone' = repoGone \\
  starts' = starts \\
  dropCause' = dropCause \\
  lostBefore' = lostN \\
  stopHeld' = outboxN + flightN \\
  stopLost' = ( \IF settled? = ztrue \THEN 0 \ELSE flightN ) \\
  lastStop' = ( \IF settled? = ztrue \THEN sGraceful \ELSE sTimedOut )
\end{schema}

\begin{schema}{ProcessLoss}
  \Delta Record \\
  \Xi Effects \\
  \Delta Mail
  \where
  recState \in \{ rsLive, rsDormant \} \\
  recState' = rsStopped \\
  currentRun' = currentRun \\
  conns' = \emptyset \\
  age' = ( \IF recState = rsLive \THEN 0 \ELSE age ) \\
  endRun' = endRun \\
  outboxN' = 0 \\
  flightN' = 0 \\
  lostN' = lostN + outboxN + flightN \\
  sessN' = sessN \\
  userN' = userN \\
  shownN' = shownN \\
  issuedN' = issuedN \\
  shownDormant' = shownDormant \\
  strandedN' = strandedN \\
  forgot' = forgot \\
  repoGone' = repoGone \\
  starts' = starts \\
  dropCause' = dropCause \\
  lostBefore' = lostN \\
  stopHeld' = outboxN + flightN \\
  stopLost' = outboxN + flightN \\
  lastStop' = sProcess
\end{schema}

A stop, of a live record or of a dormant one (a detach leaves the outbox held
with nobody connected, and the daemon may stop then): the record is durable and
is stopped, with its age kept if it was already dormant; every connection
is gone.  The outbox is the daemon's memory.  $GracefulStop$ files it back to the
session's relay inbox first, so nothing the outbox held is lost; the bounded wait
for the reads in flight can time out ($settled? = zfalse$), and the mail of a read
still in flight is then lost.  $ProcessLoss$ loses the outbox and the mail in
flight (the window the unconfirmed batch has in DES-099, PL-1753); the mail the
relay holds is not lost by either.  A stop in the middle of an end ($rsEnding$) is
not modelled.  Nor is the code's refile deadline: $GracefulStop$ files all the held mail and never
drops it at the budget, and the late broadcast fetch is a loss only through the timeout
($sTimedOut$).  Both are tracked as PL-1777 and will be modelled with its fix.

\begin{schema}{RestoreRecord}
  \Delta Record \\
  \Xi Effects \\
  \Delta Mail
  \where
  recState = rsStopped \\
  repoGone = zfalse \\
  rowPresent = ztrue \\
  recState' = rsDormant \\
  conns' = \emptyset \\
  age' = age \\
  currentRun' = currentRun \\
  endRun' = endRun \\
  outboxN' = outboxN \\
  sessN' = sessN \\
  userN' = userN \\
  shownN' = shownN \\
  issuedN' = issuedN \\
  lostN' = lostN \\
  strandedN' = strandedN \\
  flightN' = flightN \\
  shownDormant' = shownDormant \\
  forgot' = forgot \\
  lastStop' = lastStop \\
  stopHeld' = stopHeld \\
  stopLost' = stopLost \\
  lostBefore' = lostBefore \\
  repoGone' = repoGone \\
  starts' = 0 \\
  dropCause' = dropCause
\end{schema}

\begin{schema}{RestoreFail}
  \Delta Record \\
  \Xi Effects \\
  \Delta Mail
  \where
  recState = rsStopped \\
  repoGone = zfalse \\
  rowPresent = ztrue \\
  starts + 1 < maxFails \\
  recState' = rsStopped \\
  conns' = \emptyset \\
  age' = age \\
  currentRun' = currentRun \\
  endRun' = endRun \\
  outboxN' = outboxN \\
  sessN' = sessN \\
  userN' = userN \\
  shownN' = shownN \\
  issuedN' = issuedN \\
  lostN' = lostN \\
  strandedN' = strandedN \\
  flightN' = flightN \\
  shownDormant' = shownDormant \\
  forgot' = forgot \\
  lastStop' = lastStop \\
  stopHeld' = stopHeld \\
  stopLost' = stopLost \\
  lostBefore' = lostBefore \\
  repoGone' = repoGone \\
  starts' = starts + 1 \\
  dropCause' = dropCause
\end{schema}

\begin{schema}{DropRecord}
  \Delta Record \\
  \Xi Effects \\
  \Delta Mail \\
  cause? : DropCause
  \where
  recState = rsStopped \\
  ( cause? = dRepoGone \land repoGone = ztrue ) \lor
    ( cause? = dRowGone \land repoGone = zfalse \land rowPresent = zfalse ) \lor
    ( cause? = dFailedStarts \land repoGone = zfalse \land rowPresent = ztrue \land
      starts + 1 = maxFails ) \\
  recState' = rsNone \\
  conns' = \emptyset \\
  age' = 0 \\
  currentRun' = currentRun \\
  endRun' = endRun \\
  outboxN' = outboxN \\
  sessN' = 0 \\
  userN' = userN \\
  shownN' = shownN \\
  issuedN' = issuedN \\
  lostN' = lostN \\
  strandedN' = strandedN + sessN \\
  flightN' = flightN \\
  shownDormant' = shownDormant \\
  forgot' = ztrue \\
  lastStop' = lastStop \\
  stopHeld' = stopHeld \\
  stopLost' = stopLost \\
  lostBefore' = lostBefore \\
  repoGone' = repoGone \\
  starts' = 0 \\
  dropCause' = cause?
\end{schema}

\begin{schema}{HelloFirst}
  \Delta Record \\
  \Xi Effects \\
  \Delta Mail
  \where
  recState = rsStopped \\
  repoGone = zfalse \\
  recState' = rsLive \\
  conns' = \{ currentRun \} \\
  age' = 0 \\
  currentRun' = currentRun \\
  endRun' = endRun \\
  outboxN' = outboxN \\
  sessN' = sessN \\
  userN' = userN \\
  shownN' = shownN \\
  issuedN' = issuedN \\
  lostN' = lostN \\
  strandedN' = strandedN \\
  flightN' = flightN \\
  shownDormant' = shownDormant \\
  forgot' = forgot \\
  lastStop' = lastStop \\
  stopHeld' = stopHeld \\
  stopLost' = stopLost \\
  lostBefore' = lostBefore \\
  repoGone' = repoGone \\
  starts' = 0 \\
  dropCause' = dropCause
\end{schema}

\begin{schema}{RepoVanishes}
  \Xi Record \\
  \Xi Effects \\
  \Delta Mail
  \where
  repoGone = zfalse \\
  outboxN' = outboxN \\
  sessN' = sessN \\
  userN' = userN \\
  shownN' = shownN \\
  issuedN' = issuedN \\
  lostN' = lostN \\
  strandedN' = strandedN \\
  flightN' = flightN \\
  shownDormant' = shownDormant \\
  forgot' = forgot \\
  lastStop' = lastStop \\
  stopHeld' = stopHeld \\
  stopLost' = stopLost \\
  lostBefore' = lostBefore \\
  repoGone' = ztrue \\
  starts' = starts \\
  dropCause' = dropCause
\end{schema}

The daemon's records are on disk (\texttt{HarnessRecords}, one file per harness session, written
at the hello and removed by the end), so a stop leaves the record behind and a start builds the
session again.  $rsStopped$ is a record with no daemon: the process is gone, nothing is held
in memory ($outboxN = 0$, $flightN = 0$), and the relay keeps the row, the name and the mail.
At the next start each record is dealt with once:

\begin{itemize}
  \item $RestoreRecord$: the repository is there and the relay lists the session's row, so the
    session is built again, dormant, with its retention age as it was (the model's age is
    paused while the daemon is down, which is conservative: the code's clock, resumed from the
    stored $dormant\_since$, runs on).  A session that had a connection when the daemon stopped
    has no stored dormant time and starts a fresh retention, which the stops model by
    $age' = 0$ for a live record.  $starts$ goes back to $0$ ($mark$).
  \item $RestoreFail$: the start cannot build the session (the relay is down, a registration
    fails, the configuration is unusable): the record is kept and the failure counted
    ($starts$; the code retries at 5, 15 and 45 s while it runs, and the model does not
    count those, only starts).  $maxFails$ is $2$ here, $5$ in the code.
  \item $DropRecord$: the record is removed, with its cause.  The repository is gone
    ($dRepoGone$, dropped at once); the relay no longer lists the row ($dRowGone$: the session
    ended and its record outlived the end, or the row aged out while down); or the
    $maxFails$th start in a row failed ($dFailedStarts$).  The row and the tty name are not
    touched by a drop, and the mail waiting in the session's relay inbox stays there: the code
    does not store the relay in the record, so it cannot forward it.  The model counts that
    mail as $strandedN$.  This is a limit, stated and not repaired, and it is the only way mail
    is stranded: $strandedN > 0 \Rightarrow dropCause \in \{ dRepoGone, dRowGone,
    dFailedStarts \}$.
  \item $HelloFirst$: a hello for the session arrives before the start restores it, and wins;
    the failure count goes back to $0$.
  \item $RepoVanishes$: the environment; the repository directory goes.
\end{itemize}

The stop that forgets the record (the earlier model's $GracefulStopForget$ and $ProcessLossForget$, the code before the
records were kept) is not in this file: it is the control
\texttt{claude-session-end-nopersist.tex}, and it strands mail with no drop behind it.

%%----------------------------------------------------------------------
\section{Ending a session}
%%----------------------------------------------------------------------

\begin{schema}{EndOwnSession}
  \Delta Record \\
  \Delta Effects \\
  \Xi Mail \\
  run? : 0 \upto 1
  \where
  flightN = 0 \\
  recState = rsLive \\
  run? \in conns \\
  run? = currentRun \\
  recState' = rsEnding \\
  currentRun' = currentRun \\
  conns' = conns \\
  age' = age \\
  endRun' = run? \\
  rowPresent' = rowPresent \\
  ttyOwner' = ttyOwner \\
  logged' = logged \\
  rowTodo' = ztrue \\
  ttyTodo' = ztrue \\
  logTodo' = ztrue \\
  mailTodo' = ztrue \\
  newLive' = newLive \\
  ownEnded' = ztrue
\end{schema}

A connection of the current run ends its own session: the end begins and owes
its four effects.  It waits for the reads still running ($flightN = 0$), as the
code's end does.

\begin{schema}{EndRefused}
  \Xi EndState \\
  run? : 0 \upto 1
  \where
  recState = rsLive \\
  run? \in conns \\
  run? \neq currentRun
\end{schema}

A connection of a replaced run is refused and nothing changes.

\begin{schema}{ExpireClaude}
  \Delta Record \\
  \Delta Effects \\
  \Xi Mail
  \where
  flightN = 0 \\
  recState = rsDormant \\
  conns = \emptyset \\
  age = maxAge \\
  recState' = rsEnding \\
  currentRun' = currentRun \\
  conns' = conns \\
  age' = age \\
  endRun' = currentRun \\
  rowPresent' = rowPresent \\
  ttyOwner' = ttyOwner \\
  logged' = logged \\
  rowTodo' = ztrue \\
  ttyTodo' = ztrue \\
  logTodo' = ztrue \\
  mailTodo' = ztrue \\
  newLive' = newLive \\
  ownEnded' = ownEnded
\end{schema}

Retention: a dormant record with no connection, aged to the retention, begins the
same end.

\begin{schema}{DeleteRow}
  \Xi Record \\
  \Delta Effects \\
  \Xi Mail
  \where
  recState = rsEnding \\
  rowTodo = ztrue \\
  rowPresent' = zfalse \\
  rowTodo' = zfalse \\
  ttyOwner' = ttyOwner \\
  logged' = logged \\
  ttyTodo' = ttyTodo \\
  logTodo' = logTodo \\
  mailTodo' = mailTodo \\
  newLive' = newLive \\
  ownEnded' = ownEnded
\end{schema}

\begin{schema}{ReleaseTty}
  \Xi Record \\
  \Delta Effects \\
  \Xi Mail
  \where
  recState = rsEnding \\
  ttyTodo = ztrue \\
  ttyOwner' = ( \IF ttyOwner = tOld \THEN tNone \ELSE ttyOwner ) \\
  ttyTodo' = zfalse \\
  rowPresent' = rowPresent \\
  logged' = logged \\
  rowTodo' = rowTodo \\
  logTodo' = logTodo \\
  mailTodo' = mailTodo \\
  newLive' = newLive \\
  ownEnded' = ownEnded
\end{schema}

The release is owner-checked, as the code's is: it frees the name only if this
session still holds it, and a name another session took stays theirs.  The
frozen variant \texttt{claude-session-end-unordered.tex} releases by name and
frees the new session's name; \texttt{claude-session-end-unordered-owned.tex}
keeps the owner check and lets the new hello run during the end, and the
invariant holds.  So the check, not the order alone, is what protects the new
session, and the order (the client waits for the end's reply) keeps the name
from being claimed while it is still held.

\begin{schema}{AppendLogout}
  \Xi Record \\
  \Delta Effects \\
  \Xi Mail
  \where
  recState = rsEnding \\
  logTodo = ztrue \\
  logged' = ztrue \\
  logTodo' = zfalse \\
  rowPresent' = rowPresent \\
  ttyOwner' = ttyOwner \\
  rowTodo' = rowTodo \\
  ttyTodo' = ttyTodo \\
  mailTodo' = mailTodo \\
  newLive' = newLive \\
  ownEnded' = ownEnded
\end{schema}

\begin{schema}{FileOutbox}
  \Xi Record \\
  \Delta Effects \\
  \Delta Mail
  \where
  recState = rsEnding \\
  mailTodo = ztrue \\
  userN' = userN + outboxN + sessN \\
  sessN' = 0 \\
  outboxN' = 0 \\
  shownN' = shownN \\
  issuedN' = issuedN \\
  lostN' = lostN \\
  shownDormant' = shownDormant \\
  strandedN' = strandedN \\
  forgot' = forgot \\
  flightN' = flightN \\
  lastStop' = lastStop \\
  stopHeld' = stopHeld \\
  stopLost' = stopLost \\
  lostBefore' = lostBefore \\
  repoGone' = repoGone \\
  starts' = starts \\
  dropCause' = dropCause \\
  mailTodo' = zfalse \\
  rowPresent' = rowPresent \\
  ttyOwner' = ttyOwner \\
  logged' = logged \\
  rowTodo' = rowTodo \\
  ttyTodo' = ttyTodo \\
  logTodo' = logTodo \\
  newLive' = newLive \\
  ownEnded' = ownEnded
\end{schema}

The mail the session holds is filed back in the relay and the session's inbox is
drained into the user's mailbox, where any reader finds it, and the end no longer
owes it.  A mail that arrives in the session's inbox after the drain and before
the record goes stays there, as in the code (the drain stops when a fetch finds
nothing and cannot see what comes later); the model allows the arrival and does
not state a property about it.

\begin{schema}{FinishEnd}
  \Delta Record \\
  \Xi Effects \\
  \Xi Mail
  \where
  recState = rsEnding \\
  rowTodo = zfalse \\
  ttyTodo = zfalse \\
  logTodo = zfalse \\
  mailTodo = zfalse \\
  recState' = rsNone \\
  currentRun' = currentRun \\
  conns' = \emptyset \\
  age' = age \\
  endRun' = endRun
\end{schema}

The record goes only when every effect is done.

%%----------------------------------------------------------------------
\section{A clear}
%%----------------------------------------------------------------------

\begin{schema}{NewHello}
  \Xi Record \\
  \Delta Effects \\
  \Xi Mail
  \where
  recState = rsEnding \\
  newLive = zfalse \\
  newLive' = ztrue \\
  ttyOwner' = tNew \\
  rowPresent' = rowPresent \\
  logged' = logged \\
  rowTodo' = rowTodo \\
  ttyTodo' = ttyTodo \\
  logTodo' = logTodo \\
  mailTodo' = mailTodo \\
  ownEnded' = ownEnded
\end{schema}

The new session's hello after a clear.  It is enabled only when the old session's
end is complete (no record, so no end in progress) and its tty name is free, and it
takes the name.  This is the ordering the client keeps by waiting for the end's
reply before it sends the hello.

%%----------------------------------------------------------------------
\section{Mail}
%%----------------------------------------------------------------------

\begin{schema}{MailArrive}
  \Xi Record \\
  \Xi Effects \\
  \Delta Mail
  \where
  issuedN < maxMail \\
  sessN' = sessN + 1 \\
  userN' = userN \\
  issuedN' = issuedN + 1 \\
  outboxN' = outboxN \\
  shownN' = shownN \\
  lostN' = lostN \\
  shownDormant' = shownDormant \\
  strandedN' = strandedN \\
  forgot' = forgot \\
  flightN' = flightN \\
  lastStop' = lastStop \\
  stopHeld' = stopHeld \\
  stopLost' = stopLost \\
  lostBefore' = lostBefore \\
  repoGone' = repoGone \\
  starts' = starts \\
  dropCause' = dropCause
\end{schema}

A read makes two consuming fetches.  The session's own mail goes into the outbox at once
(below).  The broadcast fetch takes the user's mail, and until it returns and the read
hands the mail over to the outbox that mail is in flight, held by neither the relay nor the
outbox: the read runs on to that point whatever happens to its caller (the code's
$Uncancellable$), and an end waits for it.  A stop whose bounded wait ends first leaves
the fetch running in a process that is about to exit, and the mail it takes is lost.

\begin{schema}{UserMailArrive}
  \Xi Record \\
  \Xi Effects \\
  \Delta Mail
  \where
  issuedN < maxMail \\
  userN' = userN + 1 \\
  issuedN' = issuedN + 1 \\
  sessN' = sessN \\
  outboxN' = outboxN \\
  shownN' = shownN \\
  lostN' = lostN \\
  shownDormant' = shownDormant \\
  strandedN' = strandedN \\
  forgot' = forgot \\
  flightN' = flightN \\
  lastStop' = lastStop \\
  stopHeld' = stopHeld \\
  stopLost' = stopLost \\
  lostBefore' = lostBefore \\
  repoGone' = repoGone \\
  starts' = starts \\
  dropCause' = dropCause
\end{schema}

Mail addressed to the user, not to one session (a broadcast), arrives in the user's
mailbox.

\begin{schema}{MailFetch}
  \Xi Record \\
  \Xi Effects \\
  \Delta Mail
  \where
  recState = rsLive \\
  sessN > 0 \\
  sessN' = sessN - 1 \\
  userN' = userN \\
  outboxN' = outboxN + 1 \\
  flightN' = flightN \\
  lastStop' = lastStop \\
  stopHeld' = stopHeld \\
  stopLost' = stopLost \\
  lostBefore' = lostBefore \\
  repoGone' = repoGone \\
  starts' = starts \\
  dropCause' = dropCause \\
  shownN' = shownN \\
  issuedN' = issuedN \\
  lostN' = lostN \\
  shownDormant' = shownDormant \\
  strandedN' = strandedN \\
  forgot' = forgot
\end{schema}

The first of a read's two fetches: the session's own mail.  The model cannot represent a
crash after the relay gave the mail up and before the daemon stored it: the fetch and the
storage are one step here.  That window is PL-1753, a documented limit of this model.  The code puts it into
the outbox before the next step can run (\texttt{SessionFetch.unread} merges it at once), so
there is no moment at which it is held by neither the relay nor the outbox.

\begin{schema}{BroadcastFetch}
  \Xi Record \\
  \Xi Effects \\
  \Delta Mail
  \where
  recState = rsLive \\
  userN > 0 \\
  sessN' = sessN \\
  userN' = userN - 1 \\
  outboxN' = outboxN \\
  flightN' = flightN + 1 \\
  lastStop' = lastStop \\
  stopHeld' = stopHeld \\
  stopLost' = stopLost \\
  lostBefore' = lostBefore \\
  repoGone' = repoGone \\
  starts' = starts \\
  dropCause' = dropCause \\
  shownN' = shownN \\
  issuedN' = issuedN \\
  lostN' = lostN \\
  shownDormant' = shownDormant \\
  strandedN' = strandedN \\
  forgot' = forgot
\end{schema}

\begin{schema}{MailHandOver}
  \Xi Record \\
  \Xi Effects \\
  \Delta Mail
  \where
  recState \in \{ rsLive, rsDormant \} \\
  flightN > 0 \\
  flightN' = flightN - 1 \\
  outboxN' = outboxN + 1 \\
  sessN' = sessN \\
  userN' = userN \\
  shownN' = shownN \\
  issuedN' = issuedN \\
  lostN' = lostN \\
  shownDormant' = shownDormant \\
  strandedN' = strandedN \\
  forgot' = forgot \\
  lastStop' = lastStop \\
  stopHeld' = stopHeld \\
  stopLost' = stopLost \\
  lostBefore' = lostBefore \\
  repoGone' = repoGone \\
  starts' = starts \\
  dropCause' = dropCause
\end{schema}

The read finishes: the mail the broadcast fetch took is now held in the outbox.  The held mail is
filed back by address in the code (a broadcast message into the user's mailbox, a session
message into the session's inbox); the model counts all of it in $sessN$, which changes no
conservation count.


The client prints a message it was handed.  It needs a connected client: with the
session dormant (a daemon restart left no connection) the held mail cannot be
shown, and a reply already in flight when the daemon died counts as shown, which
is the count $shownN$ took when it was handed over.

\begin{schema}{MailShow}
  \Xi Record \\
  \Xi Effects \\
  \Delta Mail
  \where
  recState = rsLive \\
  outboxN > 0 \\
  outboxN' = outboxN - 1 \\
  shownN' = shownN + 1 \\
  sessN' = sessN \\
  userN' = userN \\
  issuedN' = issuedN \\
  lostN' = lostN \\
  shownDormant' = ( \IF recState = rsDormant \THEN ztrue \ELSE shownDormant ) \\
  strandedN' = strandedN \\
  forgot' = forgot \\
  flightN' = flightN \\
  lastStop' = lastStop \\
  stopHeld' = stopHeld \\
  stopLost' = stopLost \\
  lostBefore' = lostBefore \\
  repoGone' = repoGone \\
  starts' = starts \\
  dropCause' = dropCause
\end{schema}

\begin{schema}{Quiesce}
  \Xi EndState
\end{schema}

%%----------------------------------------------------------------------
\section{Model Checking}
%%----------------------------------------------------------------------

$EndInv$ is the state invariant.  The retention obligation is stated as
\emph{enabledness}, plus one reachability check, and not as liveness.  $EF$
alone would be satisfied by a client that reattaches and ends the session, which
is not retention, and $AF$ cannot hold of a machine whose $Quiesce$ may be taken
forever.  So:

\begin{quote}
  $AG\,(\,recState = rsDormant \land age = maxAge \Rightarrow
  enabled(ExpireClaude) \lor enabled(MailHandOver)\,)$ \\
  $AG\,(\,recState = rsDormant \land age = maxAge \Rightarrow
  EF\;enabled(ExpireClaude)\,)$ \\
  $AG\,(\,recState = rsEnding \land rowTodo = ztrue \Rightarrow
  enabled(DeleteRow)\,)$, and the same for $ReleaseTty$, $AppendLogout$ and
  $FileOutbox$ with their flags \\
  $AG\,(\,recState = rsEnding \land rowTodo = ttyTodo = logTodo = mailTodo =
  zfalse \Rightarrow enabled(FinishEnd)\,)$ \\
  $AG\,(\,recState = rsEnding \Rightarrow EF\;recState = rsNone\,)$, a
  \emph{reachability} check: from an end in progress the record can be removed
\end{quote}

$ExpireClaude$ waits for a read in flight ($flightN = 0$), as the code's
expiry does: it takes the session's lock (\texttt{held.exclusive()} in the sweep), which a
read holds from before its consuming fetch to its hand-over.  So the plain statement
$AG\,(\,recState = rsDormant \land age = maxAge \Rightarrow enabled(ExpireClaude)\,)$
is \textbf{FALSE} in this model, and says something true about the code: a dormant
record at its retention with a read still in flight is not expired until that read
hands over.  ProB's counterexample is a trace (a long one; the shortest cause is
$UserMailArrive$, $BroadcastFetch$, $Detach$, $SweepTick$, $SweepTick$) ending where $flightN = 1$ with the record dormant at the retention.
The two formulas above are what the model proves: at the retention either the expiry
or the read's hand-over is enabled, and from every such state the expiry is reachable.
The hand-over is enabled whenever $flightN > 0$ (the read runs on to it, however its caller fares).

Under the assumption that the daemon's sweep and the end's steps are scheduled
fairly (a step that stays enabled is eventually taken, and $Quiesce$ stands for the
absence of work, so it is taken only when nothing else is enabled), the
enabledness formulas give: a dormant session at its retention is removed, and an
end in progress finishes.  The model does not state fairness, and none of the
formulas is a liveness theorem.

and the non-vacuity goals are that both transitions are reachable:
$recState = rsEnding$ after $EndOwnSession$ and after $ExpireClaude$, that
$EndRefused$ is enabled ($currentRun = 1 \land 0 \in conns$), and that a clear
completes through an end ($newLive = ztrue \land ownEnded = ztrue \land forgot =
zfalse$, the ghost $ownEnded$ set by $EndOwnSession$: a trace with an
$EndOwnSession$ step before $NewHello$, not the path through the forgetting restart
and $TtyAgesOut$, which enables $NewHello$ with no end).

Each stop and each restart transition has a reachability goal, so that none is vacuous.  The stop
goals are stated over $stopHeld$ and $stopLost$, the ghosts of the last stop alone, so that mail
an earlier crash lost cannot satisfy them:
$lastStop = sGraceful \land recState = rsDormant \land forgot = zfalse \land stopHeld > 0 \land
stopLost = 0 \land sessN \geq stopHeld \land starts = 0$ (a graceful stop found mail held, in the
outbox or in flight, filed all of it back, and the record was restored: reached by
$BroadcastFetch$, $GracefulStop$, $RestoreRecord$);
$lastStop = sTimedOut \land stopLost > 0$ (the settle wait expired with a read in flight: that
read's mail, and only it, is lost);
$lastStop = sProcess \land stopLost = stopHeld \land stopLost > 0 \land recState = rsDormant$
(the crash loses all it found held, and the record comes back).
The restart goals: each cause of $DropRecord$ ($dropCause = dRepoGone$, $dRowGone$,
$dFailedStarts$, each with $recState = rsNone$); $strandedN > 0 \land dropCause = dRowGone$ (the
mail of a dropped session stays in its relay inbox); $recState = rsStopped \land starts = 1$ (a
failed start, counted); and $HelloFirst$ (a live record with $age = 0$ and $starts = 0$ after a
stop).  The enabledness formulas add
$recState = rsDormant \land outboxN > 0 \Rightarrow e(GracefulStop)$, the same for
$ProcessLoss$, $flightN > 0 \land recState \neq rsNone \Rightarrow e(MailHandOver)$, and for a
stopped record: $repoGone = zfalse \land rowPresent = ztrue \Rightarrow e(RestoreRecord)$;
$repoGone = ztrue \Rightarrow e(DropRecord)$; $repoGone = zfalse \land rowPresent = zfalse
\Rightarrow e(DropRecord)$; $starts = 1$ (the last start before the drop) with the row and the
repository there $\Rightarrow e(DropRecord) \lor e(RestoreRecord)$; $repoGone = zfalse \Rightarrow
e(HelloFirst)$.  The reachability check $AG\,(\,recState = rsStopped \Rightarrow EF\;recState \neq
rsStopped\,)$: a stopped record is always restored, dropped or reattached by a hello.

Thirteen frozen variants are the negative controls, each the shipped file with one
thing taken out:
\begin{itemize}
  \item \texttt{claude-session-end-nopersist.tex}: the records are not kept.  A stop forgets the
    record (the stops forget the record, as the earlier model's $GracefulStopForget$ and $ProcessLossForget$ did: the code before
    the records were persisted) and the mail waiting in the relay is stranded with no drop
    behind it: $strandedN > 0 \Rightarrow dropCause \in \{ dRepoGone, dRowGone, dFailedStarts \}$ is
    violated, by $MailArrive$, $ProcessLoss$ (the forgetting stop).
  \item \texttt{claude-session-end-nograceful.tex}: $GracefulStop$ skips the hand-back (the stop that does not file the held
    outbox back to the relay, the requirement that DES-101 records for a graceful
    stop): the conservation invariant is violated, by $MailArrive$, $MailFetch$,
    $GracefulStop$.
  \item \texttt{claude-session-end-noflight.tex}: $GracefulStop$
    files the outbox but not the mail of a read that settled, and do not count it as lost:
    the mail vanishes, and conservation is violated.
  \item \texttt{claude-session-end-noflight-counted.tex}: the same stops drop the settled read's
    mail and count it as lost, still reporting $sGraceful$.  Conservation holds (the loss is
    counted), and the invariant $lastStop = sGraceful \Rightarrow stopLost = 0$ is violated:
    a stop that reports success must have lost nothing.
  \item \texttt{claude-session-end-unrecorded-loss.tex}: a settled stop drops the read's mail
    and adds it to $lostN$ without reporting it in $stopLost$.  Conservation holds and
    $lastStop = sGraceful \Rightarrow stopLost = 0$ holds; $lostN = lostBefore + stopLost$ is
    violated.
  \item \texttt{claude-session-end-nofile.tex}: $FileOutbox$ drops what the session
    holds (the end that does not file mail back, the requirement A2 of DES-099
    states): the conservation invariant is violated.
  \item \texttt{claude-session-end-unordered.tex}: $NewHello$ is enabled while the
    old end is still in progress and $ReleaseTty$ frees the name by name: the old
    end frees the new session's name, and $newLive = ztrue \Rightarrow ttyOwner =
    tNew$ is violated.
  \item \texttt{claude-session-end-unordered-owned.tex}: the same unordered
    $NewHello$ with the owner-checked $ReleaseTty$ of the shipped model: the
    invariant \textbf{holds} over its 67,453 states.  This is the positive control:
    the owner check is what keeps the new session's name.
  \item \texttt{claude-session-end-orphan.tex}: $ExpireClaude$ owes no row deletion
    (the retention path forgets the row): the record goes with the row still
    present, and $rowPresent = ztrue \Rightarrow recState \neq rsNone$ is violated:
    $BadOrphan$.
  \item \texttt{claude-session-end-noexpire.tex}: the guard of $ExpireClaude$ is
    $age > maxAge$, which the saturating sweep never reaches: the first retention
    formula is \textbf{FALSE}, a dormant record at its retention is never
    removed.
  \item \texttt{claude-session-end-showanywhere.tex}: $MailShow$ needs no
    connection.  Held mail is consumed by a print with nobody connected:
    $shownDormant = zfalse$ is violated, by $UserMailArrive$, $BroadcastFetch$, $Detach$, $MailHandOver$,
    $MailShow$.  (A restart is not the route: it loses the outbox in this control
    as in the shipped model, so the held mail has nothing left to show.)
  \item \texttt{claude-session-end-anyrun.tex}: $EndOwnSession$ accepts a
    connection of a replaced run: the replaced connection ends the session the
    newer run owns (checked by the goal $recState = rsEnding$ reached with
    $run? \neq currentRun$, as the history of the operation).
\end{itemize}

%%----------------------------------------------------------------------
%% Reproducing the checks.  Run from the repo root; one probcli at a time.
%%
%%   PROBCLI="${PROBCLI:-$HOME/Applications/ProB/probcli}"
%%   P="-p DEFAULT_SETSIZE 1 -p MAXINT 6 -p MAX_INITIALISATIONS 20 \
%%      -p MAX_OPERATIONS 500 -p TIME_OUT 600000"
%%   ( ulimit -v 8000000; "$PROBCLI" docs/specs/claude-session-end.tex $P -model_check )
%%   "$PROBCLI" docs/specs/claude-session-end.tex $P -model_check -ctlformula \
%%       "AG({recState=rsDormant & age=2} => (e(ExpireClaude) or e(MailHandOver)))"
%%   (and "... => EF e(ExpireClaude)"; the plain "=> e(ExpireClaude)" is FALSE, counterexample ends BroadcastFetch, Detach, SweepTick, SweepTick)
%%   "$PROBCLI" docs/specs/claude-session-end.tex $P -model_check -ctlformula \
%%       "AG({recState=rsEnding & rowTodo=ztrue} => e(DeleteRow))"
%%   "$PROBCLI" docs/specs/claude-session-end.tex $P -model_check -ctlformula \
%%       "AG({recState=rsEnding & rowTodo=zfalse & ttyTodo=zfalse & logTodo=zfalse & mailTodo=zfalse} => e(FinishEnd))"
%%   "$PROBCLI" docs/specs/claude-session-end.tex $P -model_check -ctlformula \
%%       "AG({recState=rsEnding} => EF {recState=rsNone})"   (reachability)
%%   goals: "recState=rsEnding & age=2", "newLive=ztrue & ownEnded=ztrue & forgot=zfalse", "recState=rsLive & conns/={} & currentRun=1"
%%   "$PROBCLI" docs/specs/claude-session-end.tex $P -model_check -ctlformula \
%%       "AG({forgot=ztrue & recState=rsNone & rowPresent=ztrue} => e(RowAgesOut))"   (and TtyAgesOut)
%%   shipped model, exhaustive: 42,325 states; unordered-owned: 67,453 states, no violation
%%   goals (stops): "lastStop=sGraceful & recState=rsDormant & forgot=zfalse & stopHeld>0 & stopLost=0 & sessN>=stopHeld & starts=0",
%%       "lastStop=sTimedOut & stopLost>0", "lastStop=sProcess & stopLost=stopHeld & stopLost>0 & recState=rsDormant"
%%   goals (restart): "dropCause=dRepoGone & recState=rsNone" (and dRowGone, dFailedStarts), "strandedN>0 & dropCause=dRowGone",
%%       "recState=rsStopped & starts=1", "recState=rsLive & lastStop=sGraceful & forgot=zfalse & stopHeld>0 & age=0 & starts=0 & sessN=0"
%%   enabledness: AG({recState=rsDormant & outboxN>0} => e(GracefulStop)) (and ProcessLoss),
%%       AG({recState=rsStopped & repoGone=zfalse & rowPresent=ztrue} => e(RestoreRecord)),
%%       AG({recState=rsStopped & repoGone=ztrue} => e(DropRecord)),
%%       AG({recState=rsStopped & repoGone=zfalse & rowPresent=zfalse} => e(DropRecord)),
%%       AG({recState=rsStopped & repoGone=zfalse} => e(HelloFirst)),
%%       AG({recState=rsStopped} => EF {recState/=rsStopped}),
%%       AG({flightN>0 & recState/=rsNone} => e(MailHandOver))
%%
%% Negative controls (each reports an invariant violation, the noexpire one a
%% FALSE CTL formula):
%%
%%   for V in nopersist nograceful noflight noflight-counted unrecorded-loss nofile unordered orphan anyrun showanywhere; do
%%     "$PROBCLI" docs/specs/claude-session-end-$V.tex $P -model_check
%%   done
%%   "$PROBCLI" docs/specs/claude-session-end-noexpire.tex $P -model_check -ctlformula \
%%       "AG({recState=rsDormant & age=2} => (e(ExpireClaude) or e(MailHandOver)))"
%%----------------------------------------------------------------------

\end{document}
